% SPDX-License-Identifier: CC-BY-SA-4.0
% SPDX-FileCopyrightText: Louis Royer <infos.louis.royer@gmail.com>
\documentclass[../main.tex]{subfiles}
\begin{document}
\initSubfile
\dummylabel{ch:overview-5G-edge}{1}
\dummylabel{ch:state-of-the-art}{2}
\dummylabel{ch:edge-services-integration}{3}
\dummylabel{ch:mobility}{4}
\dummylabel{ch:testbed}{5}
\dummylabel{ch:conclusion-perspectives}{} % unnumbered

\chapter*{Introduction}
\label{ch:introduction}
\addstarredchapter{Introduction}
\markboth{Introduction}{}

\begingroup
\setlength{\parskip}{2em}

\section*{Contexte et problématique de la thèse}
Le \textit{smartphone} est désormais un équipement possédé presque universellement par la population et a su s’inscrire dans la vie quotidienne au travers d’une grande diversité d’usages~\cite{CREDOC.2026.BAROMETRE.NUMERIQUE}.
L’un des ingrédients indispensables de ce succès mondial est l’infrastructure de télécommunication qui n’a jamais cessé d’évoluer pour s’adapter aux nouvelles exigences des applications (notamment la latence et le débit).

Pour satisfaire à ces exigences et s’adapter aux besoins de chaque application, les réseaux mobiles de cinquième génération (5G) doivent pouvoir mettre en œuvre de l’\textit{edge computing} afin d’offrir des capacités de stockage et de calcul à la périphérie du réseau, proche des utilisateurs.
Le paradigme \textit{edge} permet d’améliorer la rapidité de réaction, prérequis au temps-réel, mais également la \glx{QoE}{qualité d’expérience}~globale.

Aujourd'hui, les opérateurs déploient les réseaux 5G à grande échelle, et il est attendu pour la 6G une couverture mondiale s’appuyant sur l’intégration des réseaux \glx{NTN}{non-terrestres}.
Les caractéristiques des réseaux de transports satellite, dont la bande passante est limitée et la latence élevée, amplifient le besoin d’une intégration de l’\textit{edge} aux réseaux mobiles.

Pourtant, peu de travaux de recherche ou d’implémentations \textit{open-source} ont étudié les problématiques d’intégration de l’\textit{edge} dans l’architecture~5G.
La difficulté ne réside pas seulement dans le choix de la meilleure instance \textit{edge}, mais également dans la capacité à rendre ce choix effectif dans le réseau, en configurant dynamiquement le \glx{CUPS}{plan de données} tout en restant compatible avec les mécanismes de contrôle de la~5G.


Pour permettre cette intégration, plusieurs verrous technologiques et scientifiques doivent être levés:
\begin{itemize}
    \item Rendre possible l’allocation dynamique de ressources dans l’architecture globale (dont le contrôle de leur localisation) en prenant en compte les exigences des applications;
    \item Rendre possible l’isolation des ressources de manière transparente entre les différentes communications pour pouvoir offrir des garanties de qualité de service;
    \item Gérer la mobilité afin de maintenir une continuité de service.
\end{itemize}
\endgroup

\pagebreak
Plus spécifiquement, nous étudions les alternatives d’architecture permettant l’intégration des concepts de l’\textit{edge computing} avec les réseaux~5G.
Nous cherchons à répondre à plusieurs problématiques:
\begin{itemize}
    \item Comment permettre à un opérateur d’exprimer des politiques \textit{edge} abstraites, dépendant notamment de la zone géographique et du service demandé, puis de les traduire automatiquement en règles de routage dans le \glx{CUPS}{plan de données}?
    \item Comment diriger le trafic utilisateur vers l’instance sélectionnée, de manière transparente pour l’\glx{user}{abonné} et compatible avec les réseaux~5G existants?
    \item Dans le paradigme de l’\textit{edge computing}, les services sont déployés de manière dynamique: comment prendre en compte l’évolution du déploiement?
    \item Dans les réseaux~5G, les utilisateurs sont mobiles: comment préserver les sessions existantes avec les \glx{service/instances}{services~\textit{edge}}?
    \item Comment prendre en compte les événements de mobilité pour déclencher un changement d’\glx{edge}{\textit{edge}} de manière transparente pour l’utilisateur?
\end{itemize}

\section*{Contributions}
Outre le travail d’état de l’art, nous présentons dans cette section les principales contributions qui résultent de notre travail de thèse.

\subsection*{Architecture SR4MEC pour l’intégration contrôlée de services \textit{edge} dans les réseaux 5G}
Pour répondre aux différentes problématiques énumérées dans la section précédente, nous proposons une architecture, appelée SR4MEC, qui permet à l’opérateur d’exprimer des politiques \textit{edge} de haut niveau dans le \glx{CUPS}{plan de contrôle}.
Ces politiques associent une \glx{slice}{\textit{slice}}, une zone géographique et un \glx{service/instances}{service à une instance} ou à un \glx{edge}{\textit{edge}} cible.
Elles permettent par exemple de diriger deux \glx{user}{abonnés} demandant le même \glx{service/instances}{service vers des instances} différentes selon leur région ou selon le réseau virtuel auquel appartient leur session.

L’infrastructure traduit ensuite automatiquement ces politiques abstraites en règles de configuration du \glx{CUPS}{plan de données}.
Cette traduction s’appuie sur le \glx{SR}{routage par segments}, et plus particulièrement sur \glx{SRv6}{SRv6}, afin de réaliser la redirection de trafic vers l’\gl{service/instances}{instance} spécifiée tout en réduisant les états à maintenir dans le réseau.
La sélection et l’accès à l’\glx{service/instances}{instance} sont contrôlés par l’opérateur et restent transparents pour l’\glx{UE}{équipement utilisateur}.

Cette contribution peut être décomposée comme suit:
\begin{itemize}
    \item Modélisation de politiques \textit{edge} pour mettre en relation un utilisateur avec une instance~\textit{edge};
    \item Traduction automatique des politiques en règles de routage \glx{SRv6}{SRv6} pour orienter le trafic utilisateur;
    \item Création d’une procédure de changement dynamique d’\glx{service/instances}{instance} sans interruption d’accès au \glx{service/instances}{service};
    \item Intégration de l’architecture avec les procédures de \glx{handover}{\textit{handover}}, avec préservation ou changement immédiat de l’instance~\textit{edge}.
\end{itemize}

\pagebreak
\begingroup
\setlength{\parskip}{1.5em}
\subsection*{Plateforme d’expérimentation dédiée à l’intégration de l’\textit{edge} dans les réseaux~5G}
Notre seconde contribution est la conception et le développement d’une plateforme d’expérimentation permettant d’évaluer l'intégration de \glx{service/instances}{services \textit{edge}} dans des réseaux~5G réels ou émulés.
Cette plateforme ne se limite pas à la simulation de scénarios abstraits: elle s’appuie sur des composants réseau fonctionnels, compatibles avec les interfaces du \glx{CN}{cœur~5G}, afin d'étudier concrètement l’impact de notre architecture sur le \glx{CUPS}{plan de contrôle et le plan de données}.

Un élément central de cette plateforme est l'implémentation d'un \glx{controller}{contrôleur}~SR4MEC compatible avec l’interface \glx{PFCP}{PFCP} utilisée par le \glx{SMF}{SMF}~5G.
Du point de vue du \glx{CN}{cœur de réseau}, ce \glx{controller}{contrôleur} se comporte comme un \glx{UPF}{UPF} classique.
Il reçoit donc les règles produites par le \glx{SMF}{SMF} lors de l’établissement, de la modification ou de la libération des \glx{PDU Session}{sessions~PDU}.
À partir de ces informations, il applique les politiques \textit{edge} définies par l’opérateur, puis configure automatiquement les passerelles et \textit{endpoints}~\glx{SRv6}{SRv6} afin d'orienter le trafic utilisateur vers l’\glx{service/instances}{instance de service} sélectionnée.

Cette plateforme permet ainsi de valider expérimentalement le principe de traduction entre des politiques \textit{edge} abstraites, exprimées dans le \glx{CUPS}{plan de contrôle}, et des règles concrètes de redirection de trafic dans le \glx{CUPS}{plan de données}.
Elle permet également d’évaluer les procédures de changement d’instance et de mobilité, en comparant notre architecture à une architecture~5G classique.

Cette contribution peut être décomposée comme suit:
\begin{itemize}
    \item Intégration de composants tiers pré-existants au sein d’une plateforme virtualisée et modulaire;
    \item Développement d’un \glx{UPF}{UPF}~5G classique fonctionnant dans l’espace utilisateur, servant d’implémentation de référence;
    \item Développement d’un \glx{controller}{contrôleur}~SR4MEC compatible \glx{PFCP}{PFCP}, capable d’interagir avec le \glx{SMF}{SMF}~5G, d’appliquer des politiques \textit{edge} et de configurer automatiquement le \glx{CUPS}{plan de données}~\glx{SRv6}{SRv6};
    \item Implémentation de passerelles et d’\textit{endpoints}~\glx{SRv6}{SRv6} permettant la redirection transparente du trafic utilisateur vers les \glx{service/instances}{instances \textit{edge}} sélectionnées;
    \item Développement d’un émulateur de \glx{RAN}{réseau d’accès} et de \glx{CUPS}{plan de contrôle}~5G permettant de reproduire des scénarios de mobilité et d’évaluer les procédures de \glx{handover}{\textit{handover}}.
\end{itemize}
\endgroup

\pagebreak
\section*{Plan du manuscrit}
Cette section présente la structure générale du manuscrit qui s’organise, après cette introduction, en cinq chapitres suivis d’une conclusion générale.

\begin{itemize}
    \item \hyperref[ch:overview-5G-edge]{\textbf{Chapitre~1}}:
    Dans ce chapitre, nous présentons un \textbf{aperçu des réseaux mobiles de 5\textsuperscript{ème}~génération et de l’\textit{edge computing}}
     ainsi que les enjeux liés à leur intégration.
     Nous re-traçons d’abord les évolutions des réseaux mobiles au cours des générations, puis exposons les principes de l’\textit{edge computing} comme modèle complémentaire au \textit{cloud computing} en soulignant son rôle central dans le traitement distribué de données et l’optimisation des performances.

    \item \hyperref[ch:state-of-the-art]{\textbf{Chapitre~2}}:
    Dans ce chapitre, nous effectuons un \textbf{état de l’art des travaux intégrant de l’\textit{edge} dans la 5G}.
    Nous abordons ce chapitre en définissant nos exigences pour cette intégration.
    Puis nous présentons les différents protocoles et architectures qui ont été publiés par les organismes de standardisation avant de faire la revue des travaux de recherche liés à l’\textit{edge computing}.

    \item \hyperref[ch:edge-services-integration]{\textbf{Chapitre~3}}:
    Nous présentons dans ce chapitre notre architecture permettant l’\textbf{intégration de services~\textit{edge} dans des réseaux~5G existants}.
    Nous expliquons nos choix de conception, et montrons que cette architecture permet de satisfaire aux exigences d’intégrations que nous avions définit dans le chapitre précédent.
    Nous décrivons les différentes étapes qui permettent de mettre en place et d’accéder à des services \textit{edges} dans le cas d’un utilisateur statique, et analysons l’impact de notre architecture sur le plan de contrôle.

    \item \hyperref[ch:mobility]{\textbf{Chapitre~4}}:
    Ce chapitre présente notre contribution sur la \textbf{gestion de la mobilité} et son intégration dans l’architecture que nous proposons.
    Nous montrons que notre architecture est capable de prendre en compte les \glx{handover}{\textit{handovers}} de manière compatible avec l’existant, et de préserver la continuité de la session des utilisateurs.
    Nous étudions également l’impact sur le \glx{CUPS}{plan de contrôle} des différentes procédures que nous proposons.
    \item \hyperref[ch:testbed]{\textbf{Chapitre~5}}:
    Au sein de ce chapitre, nous présentons notre \textbf{plateforme d’expérimentation} libre et open-source qui permet le déploiement d’un cœur de réseau~5G non simulé et de réseaux d’accès émulés ou réels.
    Nous utilisons cette plateforme pour \textbf{évaluer}, notre architecture la comparer à une architecture~5G classique, et la tester dans des scénarios de mobilité.

    \item \hyperref[ch:conclusion-perspectives]{\textbf{Conclusion et perspectives}}:
    Ce chapitre final fait la synthèse nos contributions et expose les perspectives d’évolution de nos travaux.

\end{itemize}
\end{document}